
%(BEGIN_QUESTION)
% Copyright 2011, Tony R. Kuphaldt, released under the Creative Commons Attribution License (v 1.0)
% This means you may do almost anything with this work of mine, so long as you give me proper credit

Du får en Fisher model 546 (eller tilsvarende) I/P-transduser for kalibrering. Du har følgende utstyr tilgjengelig:

\begin{itemize}
\item{} Justerbar DC-strømforsyning (0 til 30 volt)
\item{} Jumperledninger (krokodilleklemmer)
\item{} Multimeter med sprengt sikring (kan ikke måle strøm)
\item{} Presise trykkmålere
\item{} Presise trykkregulatorer
\item{} Mye plastrør og rørfittings
\item{} Regulert instrumentluft (25 PSI)
\item{} Kasse med elektroniske komponenter (motstander, brytere, transistorer, dioder osv.)
\item{} Grunnleggende håndverktøy (nøkler, tenger, skrutrekkere osv.)
\end{itemize}

Beskriv først hva som gjør kalibreringen vanskeligere enn normalt. Forklar deretter hvordan du kan løse utfordringen med delene du har, og skisser en enkel løsning.

\underbar{file i00719no}
%(END_QUESTION)





%(BEGIN_ANSWER)

Anbefalt poengdeling: halv uttelling for å identifisere utfordringen, halv uttelling for praktisk løsning.

\vskip 10pt

Utfordringen er at eneste måleinstrument ikke kan måle strøm (4-20 mA). En løsning er å koble en kjent motstand i serie med strømforsyningen, bruke justerbar spenning for å få variabel strøm, og beregne strøm via Ohms lov (I = V/R) ut fra målt spenning over motstanden.

%(END_ANSWER)





%(BEGIN_NOTES)

{\bf Denne oppgaven er ment for eksamen, ikke arbeidsark!}.

%(END_NOTES)
